% Quarantine-only second-revision candidate statement.
% Date: 29 August 2026.
% This file is not a live-book input and has not been Rethlas-verified.
% It makes no novelty or priority claim.  The earlier file without the
% suffixes without "revised2" are preserved as rejected provenance.

\begin{lemma}
\label{lem:p213-variation-graph-source-cutoff}
Let $K$ be a complete non-Archimedean field, possibly trivially valued,
let $Y$ be a good separated normal boundaryless $K$-analytic space that is
purely one-dimensional, and let $y\in Y^{[2,3]}$.  Let
\[
                  F_0:W_0\longrightarrow\mathbf G_{m,K}^{n_0}
\]
be a moment on an open neighborhood of $y$.  After shrinking $W_0$, one
can append a finite block $S$ of analytic units and find an open
neighborhood $W$ of $y$ and a smooth function
\[
       0\leq\chi\leq1,\qquad
       \operatorname{supp}(\chi)\Subset W,\qquad
       \chi=1\text{ on a neighborhood of }y,
\]
with the following property.  Put $F=(F_0,S)$ and
$v=F_{\mathrm{trop}}(y)$.  There is an open neighborhood
$O\subset\mathbf R^N$ of $v$ such that, for every positive
$\mathrm d'$- and $\mathrm d''$-closed ordinary Chambert--Loir--Ducros
current $T$ of bidegree $(0,0)$ on $Y$, the compactly supported
pushforward
\begin{equation}
             R_{T,F,\chi}:=(F_{\mathrm{trop}})_*(\chi T)
\label{eq:p213-source-cutoff-pushforward}
\end{equation}
is positive on $\mathbf R^N$ and is $\mathrm d'$- and
$\mathrm d''$-closed on $O$.  Its support in $O$ is contained in a finite
rational fan with vertex $v$.

Every branch of $Y$ at $y$ on which at least one coordinate of $F_0$
varies maps to a ray germ distinct from the ray germ of every other
source branch.  Consequently the coefficient of
$R_{T,F,\chi}$ on such a ray is contributed by that source branch alone.
No target-only global tube and no assertion about the exact inverse image
of a ray is made.
\end{lemma}

\begin{lemma}
\label{lem:p213-euclidean-one-dimensional-current-revised}
Let $\Omega\subset\mathbf R^N$ be open, with its standard integral affine
structure, and let $\Sigma\subset\Omega$ be a locally finite embedded
rational polyhedral complex of pure dimension one.  Let $R$ be a positive
$\mathrm d'$- and $\mathrm d''$-closed supercurrent of bidimension $(1,1)$
on $\Omega$ whose support is contained in $\Sigma$.  After subdividing
$\Sigma$, there are constants $b_e\in\mathbf R_{\geq0}$, one for every
open edge $e$, such that
\begin{equation}
                         R=\sum_e b_e[e].
\label{eq:p213-euclidean-edge-decomposition-revised}
\end{equation}
The sum is locally finite and $[e]$ is the integration current calibrated
by the primitive integral tangent of $e$.  At every vertex $v\in\Omega$
one has
\begin{equation}
                    \sum_{e\ni v}b_eu_{e,v}=0,
\label{eq:p213-euclidean-vertex-balance-revised}
\end{equation}
where $u_{e,v}$ is the primitive outgoing tangent.  In particular, $R$
has no additional current supported at a vertex.
\end{lemma}

\begin{proposition}
\label{prop:p213-normal-curve-arbitrary-valuation-current-revised}
Let $K$ be a complete non-Archimedean field, possibly trivially valued,
and let $Y$ be a connected good separated normal boundaryless
$K$-analytic space that is purely one-dimensional.  Every positive
$\mathrm d'$- and $\mathrm d''$-closed ordinary Chambert--Loir--Ducros
current $T$ of bidegree $(0,0)$ on $Y$ is a nonnegative constant multiple
of the integration current:
\begin{equation}
                         T=c\,\delta_Y,
             \qquad c\in\mathbf R_{\geq0}.
\label{eq:p213-normal-curve-arbitrary-classification-revised}
\end{equation}
Consequently $T$ is a Bedford--Taylor current in the sense of
\cite[Definition~8.6.2]{CLD25}.
\end{proposition}

\begin{proposition}
\label{prop:p213-curve-nonsmooth-locus-dichotomy-revised}
Let $K$ be a complete non-Archimedean field and let $Y$ be an irreducible
good strict reduced $K$-analytic space that is purely one-dimensional.
Exactly one of the following alternatives occurs.
\begin{enumerate}
\item The quasi-smooth locus of $Y$ is nonempty.  Its complement is a
closed locally finite set of rigid points.  If $Y$ is boundaryless, the
same assertion holds with ``rig-smooth'' in place of ``quasi-smooth''.
\item The quasi-smooth locus of $Y$ is empty.
\end{enumerate}
In particular, there is no intermediate case consisting of isolated
nonrigid failures of quasi-smoothness.  No equivalence with a statement
on the normalization is asserted.
\end{proposition}

\begin{example}
\label{ex:p213-nowhere-smooth-normal-curve-revised}
Let $K=\mathbf F_p((t))$, let $L=K(t^{1/p})$, and regard
$Y=(\mathbf A^1_L)^{\mathrm{an}}$ as a $K$-analytic curve.  Then $Y$ is
good, separated, strict, boundaryless, irreducible, normal, and purely
one-dimensional, but it is nowhere quasi-smooth over $K$.  Hence the
second alternative in
\cref{prop:p213-curve-nonsmooth-locus-dichotomy-revised} occurs, and its
nonsmooth locus contains nonrigid points.
\end{example}

\begin{remark}
\label{rem:p213-variation-graph-quarantine-scope-revised}
The preceding statements are quarantined candidates.  Their proof uses
ordinary analytic moments and ordinary CLD currents.  It does not assert
the analogous classification for weakly smooth forms presented only by
harmonic tropicalizations.  It also does not by itself settle the passage
across conductor points on a reduced nonnormal curve.
\end{remark}
